\section{Camada de Emulação}

Os recursos que não podem ser virtualizados neste trabalho são, por consequência, emulados. Como discutido anteriormente, a
emulação apresenta um custo computacional superior ao da virtualização. Em contrapartida, oferece maior versatilidade,
pois sua execução ocorre inteiramente em software.

A camada de emulação proposta baseia-se na técnica de \textit{trap-and-emulate}. Dessa forma, toda operação emulada é precedida por
uma exceção e, em algum grau, depende das características da arquitetura subjacente. Quando uma exceção ocorre, a camada de
emulação verifica se ela foi causada por uma tentativa de acesso a um recurso passível de emulação.
Em caso afirmativo, a operação correspondente é realizada em software. Caso contrário, a exceção é tratada como um erro
crítico, podendo resultar na desativação completa da \gls{vm}.

\subsection{RISC-V}

Em RISC-V, para viabilizar esse recurso, foi introduzida uma camada intermediária entre o mecanismo de tratamento de erros críticos do sistema e a exceção originalmente
gerada. Essa camada permite identificar situações passíveis de emulação antes que sejam encaminhadas ao tratador de erros críticos. Atualmente, as seguintes exceções
podem representar erros emuláveis:

\begin{itemize}
    \item \textit{Illegal Instruction}
    \item \textit{Load Access Fault}
    \item \textit{Store Access Fault}
\end{itemize}

\subsubsection{\textit{Illegal Instruction}}

A exceção \textit{Illegal Instruction} é gerada quando o processador encontra uma instrução não suportada, ou não permitida em determinado contexto.
Esse mecanismo é amplamente utilizado para emular extensões ausentes, possibilitando que a semântica da instrução seja reproduzida por software de forma transparente para o sistema em execução.
Neste trabalho, usaremos esse mecanismo para detectar leituras no \gls{csr} \textit{mtime}, alterações no estado da \gls{fpu} e para detecção de \textit{idle} que será melhor explicada na seção
de escalonamento.

\subsubsection{\textit{Load Access Fault}}
A exceção \textit{Load Access Fault} é gerada durante uma operação de leitura em uma região de memória que não pode ser acessada diretamente pelo contexto de execução atual.
No mecanismo de emulação, essa exceção permite interceptar a tentativa de leitura e fornecer os dados por meio de uma rotina de software, sendo utilizada principalmente na
emulação de dispositivos \gls{mmio}.
Quando a exceção é disparada, o hipervisor verifica inicialmente se o endereço acessado pertence a algum dispositivo \gls{mmio} associado à \gls{vm}.
Caso o endereço corresponda a um dispositivo emulado, são realizadas validações adicionais para determinar se a operação solicitada é suportada e válida para aquele dispositivo.
Somente após a conclusão dessas verificações a operação de leitura é efetivamente emulada e seu resultado retornado à \gls{vm}. O fluxo de tratamento dessa exceção no sistema
é ilustrado na Figura \ref{fig:load_access_fault_flow}.

\vspace{1cm}

\begin{figure}[h]
\centering
\begin{tikzpicture}[
    node distance=1.4cm,
    >=stealth,
    block/.style={rectangle, draw, minimum width=5cm, minimum height=1cm, align=center},
    decision/.style={diamond, draw, aspect=1, align=center},
    line/.style={draw, ->, thick}
]

\node[block, dashed] (fault) {Load Access Fault};
\node[block, below of=fault] (handler) {TrapHandler};
\node[block, below of=handler] (vcpu) {VirtualCPU};
\node[block, below of=vcpu] (vm) {VirtualMachine};
\node[block, below of=vm] (device) {Virtual Device};
\node[block, dashed, right=1cm of fault] (error) {Load Access Fault - Error};

\node[decision, right=1cm of vm] (valid) {};
\node[decision, below=1cm of device] (valid2) {};

\draw[line] (fault) -- (handler);
\draw[line] (handler) -- (vcpu);
\draw[line] (vcpu) -- (vm);
\draw[line] (vm) -- (valid);
\draw[line] (device) -- (valid2);

\draw[line] (valid) |- node[right, align=center] {Address\\Belongs\\to Device} (device);
\draw[line] (valid) -| node[right, align=center] {Address\\Does Not\\Belong\\to Any\\Device} (error.south);

\draw[line]
    (valid2.west)
	-- node[fill=white, inner sep=2pt] {Valid Operation} ++(-5,0)
    -- node[draw, dotted, fill=white,inner sep=5pt] {Read Value} ++(0,4)
    |- (handler.west);

\draw[line]
    (valid2.east)
    -- ++(1,0)
    -- node[fill=white,inner sep=2pt] {Invalid Operation} ++(8,0)
    |- (error.east);

\end{tikzpicture}
\caption{Fluxo de Tratamento de um \textit{Load Access Fault} - RISC-V.}
\label{fig:load_access_fault_flow}
\end{figure}

\subsubsection{\textit{Store Access Fault}}

A exceção \textit{Store Access Fault} é gerada durante uma operação de escrita em uma região de memória que não pode ser acessada diretamente pelo contexto de execução atual.
No mecanismo de emulação, essa exceção permite interceptar a tentativa de escrita e tratar a operação por meio de uma rotina de software, sendo
utilizada também na emulação de dispositivos \gls{mmio}.
Quando a exceção é disparada, o hipervisor verifica inicialmente se o endereço acessado pertence a algum dispositivo \gls{mmio} associado à \gls{vm}.
Caso o endereço corresponda a um dispositivo emulado, são realizadas validações adicionais para determinar se a operação solicitada é suportada e válida para aquele dispositivo.
Somente após a conclusão dessas verificações a operação de escrita é efetivamente emulada. O fluxo de tratamento dessa exceção no sistema é ilustrado na Figura \ref{fig:store_access_fault_flow}.

\vspace{1cm}

\begin{figure}[h]
\centering
\begin{tikzpicture}[
    node distance=1.4cm,
    >=stealth,
    block/.style={
        rectangle,
        draw,
        minimum width=5cm,
        minimum height=1cm,
        align=center
    },
    decision/.style={
        diamond,
        draw,
        aspect=1,
        align=center
    },
    line/.style={
        draw,
        ->,
        thick
    }
]

\node[block, dashed] (fault) {Store Access Fault};
\node[block, below=of fault] (handler) {TrapHandler};
\node[block, below=of handler] (vcpu) {VirtualCPU};
\node[block, below=of vcpu] (vm) {VirtualMachine};
\node[block, below=of vm] (device) {Virtual Device};

\node[block, dashed, right=1cm of fault] (error)
    {Store Access Fault -- Error};

\node[decision, right=1cm of vm] (valid) {};
\node[decision, below=1cm of device] (valid2) {};

	\node[block, dashed, below=2cm of valid2] (write)
    {Emulate Write};

\draw[line] (fault) -- (handler);
\draw[line] (handler) -- (vcpu);
\draw[line] (vcpu) -- (vm);
\draw[line] (vm) -- (valid);
\draw[line] (device) -- (valid2);

\draw[line]
    (valid)
    |- node[right, align=center]
       {Address\\Belongs\\to Device}
    (device);

\draw[line]
    (valid)
    -| node[right, align=center]
       {Address\\Does Not\\Belong\\to Any\\Device}
    (error.south);

\draw[line]
    (valid2) -- ++(0,-2)
    node[midway, fill=white, inner sep=2pt] {Valid Operation}
    -- (write);

\draw[line]
    (valid2.east)
    -- ++(1,0)
    -- node[fill=white, inner sep=2pt]
       {Invalid Operation}
    ++(8,0)
    |- (error.east);

\end{tikzpicture}
\caption{Fluxo de Tratamento de um \textit{Store Access Fault} -- RISC-V.}
\label{fig:store_access_fault_flow}
\end{figure}
